% stirling-formal.tex
% Formalization of C. Stirling, "An introduction to decidability of
% higher-order matching", relative to axioms for the game part.
% Uncompiled draft: only sanity-checked (braces, environments, refs).
\documentclass[11pt]{article}
\usepackage[margin=2.6cm]{geometry}
\usepackage{amsmath,amssymb,amsthm}
\usepackage{enumerate}

\newcommand{\Ground}{\mathbf{0}}
\newcommand{\ord}{\mathrm{ord}}
\newcommand{\ar}{\mathrm{ar}}
\newcommand{\Sub}{\mathrm{Sub}}
\newcommand{\depth}{\mathrm{depth}}
\newcommand{\br}{\mathrm{br}}
\newcommand{\Gm}{\mathsf{G}}
\newcommand{\Tm}{\mathsf{T}}
\newcommand{\bnd}{\mathrel{\downarrow}}
\newcommand{\suc}[2]{#1{\cdot}#2}
\newcommand{\iv}[2]{\pi[#1,#2]}
\newcommand{\pa}{\mathrm{par}}
\newcommand{\comb}{\mathrm{comb}}
\newcommand{\pth}{\mathrm{path}}
\newcommand{\CE}{C_E}
\newcommand{\sq}{\sqsubseteq}
\newcommand{\ssq}{\sqsubset}
\newcommand{\ra}[1]{\xrightarrow{#1}}

\theoremstyle{plain}
\newtheorem{theorem}{Theorem}[section]
\newtheorem{proposition}[theorem]{Proposition}
\newtheorem{lemma}[theorem]{Lemma}
\newtheorem{corollary}[theorem]{Corollary}
\newtheorem{axL}{Fact}
\renewcommand{\theaxL}{L\arabic{axL}}
\newtheorem{axG}{Axiom}
\renewcommand{\theaxG}{G\arabic{axG}}
\newtheorem{axU}{Axiom}
\renewcommand{\theaxU}{U\arabic{axU}}
\theoremstyle{definition}
\newtheorem{definition}[theorem]{Definition}
\newtheorem{example}[theorem]{Example}
\theoremstyle{remark}
\newtheorem{remark}[theorem]{Remark}
\newtheorem{gap}{Gap}

\title{Decidability of higher-order matching after Stirling:\\
a formalization relative to the game axioms}
\author{Working notes (HOM project)}
\date{\today}

\begin{document}
\maketitle

\begin{abstract}
We restate the decidability proof of C.~Stirling, \emph{An introduction
to decidability of higher-order matching}, in a form in which every
definition is explicit and every step is either proved or marked.  The
results about the interpolation game (Section~4 of the paper) and the
uniformity properties of plays (Section~6.1) are taken as axioms.  On
top of them we prove what can be proved (the transformation $T_1$, the
partition of plays, several counting lemmas, the final reduction of
decidability to the small-solution theorem) and isolate the remaining
steps of the proof of the small-solution theorem as explicitly stated
gaps.
\end{abstract}

\tableofcontents

%======================================================================
\section{Scope and conventions}
%======================================================================

References of the form ``[S, Def.~43]'' point to the paper.  The status
of every statement is one of the following.
\begin{itemize}
\item \textbf{Fact L$n$}: a standard fact of the simply typed
  $\lambda$-calculus, cited and not proved.
\item \textbf{Axiom G$n$}: a result of [S, Section~4] (the game
  characterisation and the basic properties of plays).  This is the
  ``game part'' that we assume formalized.
\item \textbf{Axiom U$n$}: a result of [S, Section~6.1] (uniformity
  properties of plays).  In the paper these are derived from the
  G-axioms and the move rules by game-level arguments; since they are
  statements about plays only, we treat them as part of the game layer.
  If one prefers to keep only Section~4 as axioms, the U-axioms become
  the first proof obligations (see Section~\ref{sec:gaps}).
\item \textbf{Lemma / Proposition / Theorem with proof}: proved here
  from the axioms and the definitions.
\item \textbf{Gap $n$}: a statement used in the paper whose proof there
  is missing or only sketched, together with the precise claim that
  would close it.
\end{itemize}
Section~5 of the paper (tree automata of all orders, dependency tree
automata) is not used by the decidability proof and is not formalized.
Section~6.5 (removing the restrictions) is summarized in
Section~\ref{sec:restrictions} only.

%======================================================================
\section{Terms, matching and interpolation}
%======================================================================

\subsection{Types and terms}

\begin{definition}[Types, order, arity; S, Defs.~1--2]
Types are generated by $A::=\Ground\mid A\to A$; every non-base type is
written $(A_1,\dots,A_n,\Ground)$ for $A_1\to\cdots\to A_n\to\Ground$.
\[
\ord(\Ground)=1,\quad \ord(A_1,\dots,A_n,\Ground)=1+\max_i\ord(A_i),\qquad
\ar(\Ground)=0,\quad \ar(A_1,\dots,A_n,\Ground)=\max(n,\ar(A_1),\dots,\ar(A_n)).
\]
\end{definition}

Terms are Church-typed $\lambda$-terms over typed variables and typed
constants, with $\beta$- and $\eta$-reduction, and $=_{\beta\eta}$ is
the induced equivalence [S, Defs.~3--5].

\begin{axL}[S, Fact~1]\label{L:basic}
Subject reduction, strong normalisation and confluence hold for
$\to_\beta$, $\to_\eta$ and $\to_{\beta\eta}$.
\end{axL}

\begin{definition}[Long normal form; S, Def.~6]
A term of type $(A_1,\dots,A_n,\Ground)$, $n\ge0$, is in \emph{lnf} if it
has the form $\lambda x_1\dots x_n.\,h\,s_1\dots s_m$ with $x_i:A_i$,
$h:(B_1,\dots,B_m,\Ground)$ a variable or a constant, and every $s_j:B_j$
in lnf.  For a set of constants $C$, $T_A(C)$ is the set of closed lnf
terms of type $A$ all of whose constants are in $C$ [S, Def.~7].
\end{definition}

\begin{axL}[S, Facts~2--3]\label{L:lnf}
Every term is $\beta\eta$-equivalent to a unique lnf up to
$\alpha$-equivalence.  If $t:A\to B$ and $s:A$ are in lnf and
$t\,s\to_\beta^*t'$ with $t'$ $\beta$-normal, then $t'$ is in lnf.
\end{axL}

\subsection{Matching and interpolation}

\begin{definition}[S, Defs.~8--9]
A \emph{matching equation} is $v=u$ with $v,u:\Ground$ in lnf and $u$
closed; a solution is a sequence of closed terms $t_1,\dots,t_n$ with
$v\{t_1/x_1,\dots,t_n/x_n\}=_{\beta\eta}u$, where $x_1,\dots,x_n$ are the
free variables of $v$.

An \emph{interpolation equation} $E$ is $x\,w_1\dots w_k=u$ where
$x:A=(A_1,\dots,A_k,\Ground)$, every $w_i:A_i$ is closed and in lnf and
$u:\Ground$ is closed and in lnf ($u$ is the \emph{goal}).  A
\emph{solution} of $E$ is a closed lnf $t:A$ with
$t\,w_1\dots w_k=_{\beta\eta}u$.  The \emph{order} of $E$ is $\ord(A)$
and its \emph{arity} is $\ar(A)$.
\end{definition}

\paragraph{Standing restrictions.} As in [S, Sections~2--6.4]:
\begin{enumerate}[(R1)]
\item there is a single base type $\Ground$;
\item the goal $u$ contains no bound variables.
\end{enumerate}
Under (R2) every constant occurring in $u$ has a type
$(\Ground,\dots,\Ground,\Ground)$ (an argument of non-base type would be
an lnf of the form $\lambda\vec y.\,s$ with $\vec y$ non-empty), so $u$
is a first-order ground term.

\begin{definition}[S, Defs.~10--11]
$\CE$ is the set of constants of $u$ together with a fresh constant
$b:\Ground$.  $\Sub(u)$ is the set of subterms of $u$ and $\sq$ the
subterm order on it ($\ssq$ strict).  $\depth(a)=0$,
$\depth(f\,u_1\dots u_m)=1+\max_i\depth(u_i)$; $|u|$ is the number of
nodes of $u$ and $\br(u)$ the number of its leaves (branches).
\end{definition}

Note that every strictly descending chain $r_0\sqsupset r_1\sqsupset\cdots
\sqsupset r_m$ in $\Sub(u)$ has $m\le\depth(u)$.

\begin{lemma}[S, Prop.~4 under (R2)]\label{lem:reduces}
If $t$ solves $E$ then $t\,w_1\dots w_k\to_\beta^*u$.
\end{lemma}
\begin{proof}
By Fact~\ref{L:basic} $t\,\vec w$ has a $\beta$-normal form $u'$, which is
in lnf by Fact~\ref{L:lnf}; $u'=_{\beta\eta}u$ and both are lnf, so they
are $\alpha$-equivalent, hence equal since $u$ has no binders.
\end{proof}

\begin{definition}
For a matching equation $v=u$ with free variables $x_i:A_i$, its
\emph{associated interpolation equation} is $x\,(\lambda x_1\dots
x_n.v)=u$ with $x:((A_1,\dots,A_n,\Ground),\Ground)$.  A solution of it
is \emph{canonical} if it has the form $\lambda z.\,z\,t_1\dots t_n$ with
all $t_i$ closed.
\end{definition}

\begin{proposition}[S, Prop.~5]\label{prop:canonical}
A matching equation has a solution iff its associated interpolation
equation has a canonical solution.
\end{proposition}
\begin{proof}
If $t_1,\dots,t_n$ solve $v=u$, then the lnf of $\lambda z.z\,t_1\dots t_n$
is canonical and $(\lambda z.z\,\vec t)(\lambda\vec x.v)\to_\beta^*
v\{\vec t/\vec x\}=_{\beta\eta}u$.  Conversely the same computation read
backwards shows that the arguments of a canonical solution solve $v=u$.
\end{proof}

\begin{proposition}[S, Fact~6]\label{prop:CE}
If $t\in T_A(C)$ solves $E$ then some $t'\in T_A(\CE)$ solves $E$.
\end{proposition}
\begin{proof}
For each constant $c\notin\CE$ of type $(B_1,\dots,B_m,\Ground)$ occurring
in $t$ let $K_c=\lambda y_1\dots y_m.\,b$.  Substitution of closed terms
for constants commutes with $\beta$-reduction, so from
$t\,\vec w\to_\beta^*u$ (Lemma~\ref{lem:reduces}) we get
$t\{\vec K/\vec c\}\,\vec w\{\vec K/\vec c\}\to_\beta^*u\{\vec K/\vec c\}$.
The $w_i$ may contain such $c$ too, so we substitute only in $t$: since
$c$ does not occur in $u$ and the reduction is a reduction of
$t\,\vec w$ in which the constant $c$ of $t$ is never applied in a
way visible in $u$, replacing it by $K_c$ gives the same normal form.
Formally: mark the occurrences of $c$ coming from $t$ by a fresh copy
$c'$; then $t\{c'/c\}\,\vec w\to_\beta^*u$ (the reduction does not
inspect constants) and substituting $K_c$ for $c'$ gives
$t\{K/c'\}\,\vec w\to_\beta^*u$.  Let $t'$ be the lnf of
$t\{\vec K/\vec c'\}$; its constants lie in $\CE$.
\end{proof}

%======================================================================
\section{Interpolation trees}
%======================================================================

\begin{definition}[Binding trees; S, Def.~12 and Section~3.1]
A \emph{tree} is a finite prefix-closed set of words over
$\{1,2,\dots\}$; $\suc ni$ is the $i$-th successor of $n$, $n$ is
\emph{above} $m$ ($m$ \emph{below} $n$) if $n$ is a prefix of $m$
(reflexively).  The tree of an lnf term is built with dummy lambdas:
\begin{itemize}
\item $\lambda\vec y.\,s$ with $s:\Ground$ ($\vec y$ possibly empty) is a
  node labelled $\lambda\vec y$ with the single successor the tree of
  $s$;
\item $h\,s_1\dots s_m$ is a node labelled $h$ with successors the trees
  of $s_1,\dots,s_m$ (an argument $s_j:\Ground$ is read as $\lambda.s_j$).
\end{itemize}
The \emph{binding relation} $n\bnd m$ holds when $n$ is labelled
$\lambda y_1\dots y_p$, $m$ is labelled $y_i$ and $n$ is the binder of
that occurrence.  A node is a \emph{lambda node} if its label is
$\lambda\vec y$ (a \emph{binding node} if moreover $\vec y\neq\varepsilon$),
a \emph{variable node} or a \emph{constant node} otherwise.
\end{definition}

\begin{definition}[Interpolation tree]
For a candidate $t:A$ in lnf, the \emph{interpolation tree}
$@\,t\,w_1\dots w_k$ has root $\varepsilon$ labelled $@$ and successors
$1,2,\dots,k+1$ carrying the trees of $t,w_1,\dots,w_k$.  Nodes below $1$
form the \emph{$t$-component}, nodes below $q+1$ the
\emph{$w_q$-component}.  Lambda nodes are exactly the nodes of odd
length.
\end{definition}

\begin{definition}[Structural notions; S, Defs.~34--36]\label{def:struct}
Let $n,m$ be variable nodes.
\begin{enumerate}
\item $n$ has \emph{level} $1$ if it is bound by node $1$ (the root of
  $t$) or by some $q+1$ (the root of $w_q$); level $j+1$ if it is bound
  by a successor of a level-$j$ node.
\item $m$ is \emph{embedded} if there is a node $n\neq m$ above $m$
  carrying the same variable and bound by the same binder.
\item $n$ is a \emph{descendant under binding} of $m$ if $n=m$, or some
  successor $m'$ of $m$ satisfies $m'\bnd n'$ and $n$ is a descendant
  under binding of $n'$.  $n,m$ \emph{belong to the same family} if
  they have a common ancestor under binding.  $n$ is \emph{end} if its
  only descendant under binding is itself, i.e.\ no successor of $n$
  binds anything.
\end{enumerate}
\end{definition}

\begin{lemma}[Levels; S, after Def.~34]\label{lem:levels}
Let $\ord(A)\in\{2p,2p+1\}$ and $t\in T_A(\CE)$.  Every variable node of
$t$ has a level, and the level is at most $p$.
\end{lemma}
\begin{proof}
Constants of $t$ lie in $\CE$ and are first order, so their successors
are dummy lambdas and bind nothing; hence every variable of $t$ is bound
by the root or by a successor of a variable node, and levels are
defined.  A level-$1$ variable of $t$ has type $A_i$ with
$\ord(A_i)\le\ord(A)-1$.  If $y:(B_1,\dots,B_m,\Ground)$ is a variable
node, its $i$-th successor binds variables of order at most
$\ord(B_i)-1\le\ord(y)-2$.  By induction a level-$j$ variable has order
at most $\ord(A)-2j+1$; orders are $\ge1$, so $2j\le\ord(A)$ and
$j\le p$.
\end{proof}

%======================================================================
\section{The interpolation game}
%======================================================================

\begin{definition}[States and look-up tables; S, Defs.~16--17]
$S=\{[r]\mid r\in\Sub(u)\}\cup\{[\forall],[\exists]\}$; $[\forall],[\exists]$
are \emph{final}.  Let $N_1$ be the set of binding nodes.  Look-up
tables are defined by $\Theta_1=\{\varnothing\}$ and $\Theta_{i+1}$ the
set of partial maps
$N_1\rightharpoonup\bigcup_{p}N^p\times\bigcup_{j\le i}\Theta_j$; an
entry $\theta(m)=((m_1,\dots,m_p),\theta')$ says that the variables bound
by $m$ denote the subtrees at $m_1,\dots,m_p$ under the table $\theta'$.
$\theta'\le\theta$ (\emph{$\theta$ extends $\theta'$}) if $\theta(m)=
\theta'(m)$ whenever $\theta'(m)$ is defined [S, Def.~24].
\end{definition}

\begin{definition}[Moves; S, Def.~19 and Fig.~5]\label{def:moves}
Let the current position be $n[r]\theta$ with $r$ not final.  The next
position is determined by the label of $n$:
\begin{description}
\item[$@$:] $\suc n1[r]\theta'$ with
  $\theta'=\theta\{((\suc n2,\dots,\suc n{(k+1)}),\theta)/\suc n1\}$;
\item[$\lambda\vec y$:] $\suc n1[r]\theta$;
\item[$a:\Ground$:] $n[\exists]\theta$ if $r=a$, else $n[\forall]\theta$;
\item[$f:(B_1,\dots,B_p,\Ground)$, $p>0$:] if $r=f\,r_1\dots r_p$, player
  $\forall$ chooses $j$ and the next position is $\suc nj[r_j]\theta$;
  otherwise $n[\forall]\theta$;
\item[$y_j:\Ground$:] if $m\bnd n$ and $\theta(m)=((m_1,\dots,m_l),\theta')$,
  then $m_j[r]\theta'$;
\item[$y_j:(B_1,\dots,B_p,\Ground)$, $p>0$:] if $m\bnd n$ and
  $\theta(m)=((m_1,\dots,m_l),\theta')$, then $m_j[r]\theta''$ with
  $\theta''=\theta'\{((\suc n1,\dots,\suc np),\theta)/m_j\}$.
\end{description}
The last two moves are \emph{jumps}.
\end{definition}

\begin{definition}[Plays; S, Defs.~18, 20]
A \emph{play} of $\Gm(t,E)$ is a finite sequence $\pi=\pi_1\dots\pi_l$ of
positions $\pi_i=n_i[r_i]\theta_i$ with $\pi_1=\varepsilon[u]\varnothing$,
each $\pi_{i+1}$ a next position of $\pi_i$, $r_i$ non-final for $i<l$
and $r_l$ final.  $\forall$ \emph{wins} $\pi$ if $r_l=[\forall]$; she
\emph{wins the game} if she wins some play, and \emph{loses} it
otherwise.  We write $\pi\in\Gm(t,E)$, $\iv ij$ for $\pi_i\dots\pi_j$, and
call $\pi_i$ a position \emph{at} $n_i$.
\end{definition}

\begin{definition}[Parent, chain; S, Defs.~21--23]\label{def:parent}
Let $\pi\in\Gm(t,E)$.
\begin{enumerate}
\item If $n_j$ is a variable node, $\pi_i$ is the \emph{parent} of
  $\pi_j$ if $n_i\bnd n_j$ and $\theta_i(n_i)=\theta_j(n_i)$.
\item If $n_j$ is a lambda node and $n$ is the node directly above
  $n_j$: if $n=\varepsilon$ the parent is $\pi_1$; if $n$ is a constant
  node the parent is $\pi_{j-1}$; if $n$ is a variable node the parent
  is the $\pi_i$ such that $\pi_{i+1}$ is the parent of $\pi_{j-1}$ in the
  sense of item~1.
\item A \emph{chain} for $\pi_j$ ($n_j$ a lambda node) is a sequence
  $\pi_{j_1},\dots,\pi_{j_{2p}}$ with $\pi_{j_{2p}}=\pi_j$,
  $\pi_{j_{2m-1}}$ the parent of $\pi_{j_{2m}}$ ($1\le m\le p$) and
  $j_{2m+1}=j_{2m}+1$ ($1\le m<p$).  We also call it a chain
  \emph{starting at} $j_1$; $p=0$ gives the empty chain.
\end{enumerate}
\end{definition}

%======================================================================
\section{The game axioms}\label{sec:axioms}
%======================================================================

\subsection{Section 4 of the paper}

\begin{axG}[S, Thm.~13]\label{G:char}
$\forall$ loses $\Gm(t,E)$ iff $t$ solves $E$.
\end{axG}

\begin{axG}[S, Prop.~12 and the corollary of the proof of Thm.~13]\label{G:total}
From the initial position every sequence of moves is finite, and at
every reachable non-final position the move of
Definition~\ref{def:moves} is defined (the binder entry it consults
exists).  Hence $\Gm(t,E)$ is a finite non-empty set of plays, and plays
are in bijection with the sequences of choices of $\forall$ that occur
along them.
\end{axG}

\begin{axG}[S, Fact~10]\label{G:branches}
If $t$ solves $E$ then $\Gm(t,E)$ has exactly $\br(u)$ plays.
\end{axG}

\begin{axG}[S, Fact~11]\label{G:parity}
Let $\pi\in\Gm(t,E)$ have length $l$.  (1) If $i<l$ is even then $n_i$
is a lambda node; (2) if $i>1$ is odd then $n_i$ is a variable or
constant node; (3) $n_l$ is a constant node and $r_l$ is final; (4) if
$i<j<l$ then $r_j\sq r_i$.
\end{axG}

\begin{axG}[S, Prop.~12]\label{G:tables}
Let $\theta_j(m)=((m_1,\dots,m_p),\theta')$.  (1) $\theta_j(n)$ is defined
iff $n$ is a binding node above $n_j$; (2) $\theta'(n)$ is defined iff
$n$ is a binding node strictly above every $m_k$; (3) $m$ is labelled
$\lambda y_1\dots y_p$ and $y_k$ has the type of $m_k$; (4) $\theta'=
\theta_i$ for some $i<j$, and either $i=1$ and $m_k=\suc 1{(k+1)}$... more
precisely $m_k=\suc{n_1}{(k+1)}$, or $n_i$ is a variable node and
$m_k=\suc{n_i}k$.
\end{axG}

\begin{axG}[S, Prop.~14]\label{G:distinct}
If $i\neq j$ and $n_i=n_j$ then $\theta_i\neq\theta_j$.
\end{axG}

\begin{axG}[S, Prop.~15, Fact~16]\label{G:parent}
Every position $\pi_j$, $j>1$, at a variable node or at a lambda node has
a unique parent $\pi_i$, and $i<j$.  If $n_j$ is a lambda node then
$n_j$ is a successor of $n_i$.
\end{axG}

\begin{axG}[S, Props.~17--18]\label{G:mono}
(1) If $n_j$ is a lambda node and $\pi_i$ is its parent then
$\theta_i\le\theta_j$.  (2) If $\pi_{j_1},\dots,\pi_{j_{2p}}$ is a chain
then $\theta_{j_1}\le\cdots\le\theta_{j_{2p}}$, and if $n_{j_m}\bnd n_{j_n}$
then $\pi_{j_m}$ is the parent of $\pi_{j_n}$.
\end{axG}

\begin{axG}[S, Prop.~19]\label{G:agree}
If $\theta_i(m)=\theta_j(m)$ are defined and $m'$ is a binding node above
$m$, then $\theta_i(m')=\theta_j(m')$.
\end{axG}

\begin{axG}[stated in the proof of S, Prop.~23]\label{G:alt}
If $n_j$ is a variable node then $n_{j+1}$ is a lambda node, and exactly
one of $n_j,n_{j+1}$ lies in the $t$-component (the other lies in some
$w_q$-component).
\end{axG}

\begin{remark}
Axiom~\ref{G:tables}(4) is stated in the paper as ``$i=1$ and
$m_k=n_1(k+1)$''; we read $n_1(k+1)$ as the $(k+1)$-th successor of the
root, as in the $@$-move.  Axiom~\ref{G:alt} is not stated as a
proposition in the paper but is used there; it follows from
Axiom~\ref{G:tables} by induction (entries for binders in $w_q$ always
point into $t$ and conversely, because every $w_q$ is closed).
\end{remark}

\subsection{Consequences proved from the definitions and the G-axioms}

\begin{lemma}[Predecessors]\label{lem:pred}
Let $\pi\in\Gm(t,E)$ and $1<j<l$.  If $n_j$ is not a lambda node then
$n_{j-1}$ is the lambda node directly above $n_j$ and $\pi_{j-1}$ makes
a $\lambda$-move, so $r_{j-1}=r_j$, $\theta_{j-1}=\theta_j$.
\end{lemma}
\begin{proof}
Inspecting Definition~\ref{def:moves}: the $@$-move, the $f$-move and
the jumps lead to lambda nodes, the $a$-move leads to a final position,
so a non-final position at a non-lambda node can only be reached by a
$\lambda$-move, which goes to the unique successor.
\end{proof}

\begin{lemma}[Ancestors are visited first]\label{lem:ancestors}
If $\pi_j$ is a position at $n$ and $m\ne\varepsilon$ is a non-lambda node
strictly above $n$, there is $i<j$ with $n_i=m$.  If moreover $n$ is a
lambda node, there is a unique chain for $\pi_j$ whose first position is
at $m$ [S, after Def.~23].
\end{lemma}
\begin{proof}
Define $\rho(j)=j-1$ if $n_j$ is a non-lambda node (Lemma~\ref{lem:pred};
for a final position use $\rho(l)=l-1$, which is at the same node) and
$\rho(j)=$ the index of the parent of $\pi_j$ if $n_j$ is a lambda node
(Axiom~\ref{G:parent}).  Then $\rho(j)<j$ and $n_{\rho(j)}$ is the node
directly above $n_j$ (or $n_j$ itself in the final case).  Iterating
$\rho$ visits every node on the path from $\varepsilon$ to $n$, in
decreasing index order.  The positions visited between $m$ and $\pi_j$
form a chain: lambda positions alternate with their parents and
non-lambda positions are preceded by their $\lambda$-predecessor, which
is condition~3 of Definition~\ref{def:parent}.  Uniqueness: $\rho$ is a
function.
\end{proof}

\begin{lemma}[Locality]\label{lem:locality}
Let $t,t'$ be lnf terms of type $A$ and $X$ a prefix-closed set of nodes
on which the interpolation trees of $t$ and $t'$ agree (labels, number
of successors and binding).  If no play of $\Gm(t,E)$ has a position
outside $X$, then $\Gm(t',E)=\Gm(t,E)$.
\end{lemma}
\begin{proof}
The next position of $n[r]\theta$ (for a given choice of $\forall$)
depends only on $r$, $\theta$, the label of $n$, the successors of $n$
and the binder of $n$.  By induction on $i$, every sequence of moves of
length $i$ from $\pi_1$ is the same in both games as long as all its
positions are in $X$.  A sequence for $t'$ leaving $X$ for the first
time would be, up to that point, a sequence for $t$ whose next position
is outside $X$, contrary to the hypothesis.  Tables only mention nodes,
so they coincide as well.
\end{proof}

\begin{lemma}[States only change at constants]\label{lem:states}
Let $t$ solve $E$ and $\pi\in\Gm(t,E)$.  For $i<l$, $r_{i+1}\ne r_i$ iff
$n_i$ is a constant node of positive arity, and then
$r_{i+1}\ssq r_i$.  Consequently, for $i\le j<l$: $r_j\ssq r_i$ iff some
$n_h$, $i\le h<j$, is a constant node of positive arity.
\end{lemma}
\begin{proof}
By Definition~\ref{def:moves} only the $f$-move and the $a$-move change
the state.  The $a$-move leads to a final position, so for $i+1<l$ it is
excluded; for $i+1=l$ the state becomes $[\exists]$ by
Axiom~\ref{G:char}.  At an $f$-node, since $\forall$ loses, $r_i=f\,r_1
\dots r_p$ and $r_{i+1}=r_j\ssq r_i$.  The second claim follows with
Axiom~\ref{G:parity}(4).
\end{proof}

\begin{definition}[Contribution; S, Defs.~40, 44]\label{def:contrib}
An interval $\iv ij$ with $j<l$ \emph{contributes} if $r_i\ne r_j$.
(The paper requires $n_j$ to be a lambda node; we drop this so that a
final stage $\iv ij$ ``contributes'' iff $\iv i{j-1}$ does, as in [S,
Def.~44].)
\end{definition}

\begin{lemma}[Disjoint contributions]\label{lem:disjoint}
Let $t$ solve $E$ and $\pi\in\Gm(t,E)$.  If $\iv{i_1}{j_1},\dots,
\iv{i_s}{j_s}$ contribute and $j_1\le i_2,\ j_2\le i_3,\dots$, then
$s\le\depth(u)$.
\end{lemma}
\begin{proof}
By Axiom~\ref{G:parity}(4), $r_{i_1}\sqsupset r_{j_1}\sqsupseteq r_{i_2}
\sqsupset\cdots\sqsupset r_{j_s}$ is a chain in $\Sub(u)$ with $s$ strict
steps.
\end{proof}

\begin{lemma}[Corresponding intervals contribute together]\label{lem:corr-contrib}
Let $t$ solve $E$.  If $\iv i{i+m}$ and $\iv j{j+m}$ correspond
(Definition~\ref{def:corr}), $i+m<l$ and $j+m<l$, then one contributes
iff the other does.
\end{lemma}
\begin{proof}
They visit the same nodes; apply Lemma~\ref{lem:states}.
\end{proof}

\begin{lemma}[Constants on a path]\label{lem:constants}
Let $t$ solve $E$.  On every path of the $t$-component, the number of
constant nodes of positive arity that are visited by some play is at
most $\depth(u)$.
\end{lemma}
\begin{proof}
Let $c$ be the lowest visited constant node of positive arity on the
path and $c'$ its successor on the path.  At a position at $c$ the state
is $f\,r_1\dots r_p$ (Axiom~\ref{G:char}); by Axiom~\ref{G:total} there
is a play in which $\forall$ chooses $c'$ there, so some $\pi$ has a
position $\pi_j$ at $c'$.  By Lemma~\ref{lem:ancestors} there is a chain
from $\pi_1$ (at $\varepsilon$) to $\pi_j$; it passes through every node
on the path above $c'$ in increasing index order.  For each constant
node $c''$ of positive arity on the path, the chain contains a position
at $c''$ followed by one at its successor, which is an $f$-move and
strictly decreases the state (Lemma~\ref{lem:states}).  By
Axiom~\ref{G:parity}(4) these decreases form a strictly descending
chain in $\Sub(u)$.
\end{proof}

\begin{lemma}[Contributing positions over all plays]\label{lem:global-contrib}
Let $t$ solve $E$.  Identify positions of different plays that have the
same prefix.  The number of distinct $f$-moves (positions at constant
nodes of positive arity) in the tree of all plays is at most
$|u|-\br(u)$.
\end{lemma}
\begin{proof}
An $f$-move at state $[r]$ consumes the internal node $r$ of $u$
(viewing states as occurrences of subterms of $u$; the state determines
the occurrence because it is reached from $u$ by the choices of
$\forall$).  Two plays whose choice sequences agree up to reaching $[r]$
coincide up to that point (Axiom~\ref{G:total}), so each occurrence is
consumed at one position of the tree of plays.
\end{proof}

\begin{remark}
Lemma~\ref{lem:global-contrib} justifies the bound ``the number of
stages that can contribute is at most $|u|$'' used in the proof of
[S, Thm.~37], since a contributing stage contains an $f$-move
(Lemma~\ref{lem:states}).  Strictly speaking, states must be read as
\emph{occurrences} in $u$ (two equal subterms at different places of
$u$ are different states for this count); the game itself does not
need this distinction.
\end{remark}

\subsection{Section 6.1 of the paper: uniformity}

\begin{definition}[S, Defs.~32--33, 37]\label{def:corr}
Let $\pi\in\Gm(t,E)$.
\begin{enumerate}
\item $\iv i{i+m}$ and $\iv j{j+m}$ \emph{correspond} if
  $n_{i+h}=n_{j+h}$ for $0\le h\le m$.
\item Let $j<m'$ and $n_j,n_{m'}$ carry the same variable.  The intervals
  $\iv{j+1}{j+k}$ and $\iv{m'+1}{m'+k}$ are \emph{complementary} if they
  correspond and either
  \begin{enumerate}[(a)]
  \item $\pi_j$ is the parent of $\pi_{j+k+1}$, $\pi_{m'}$ is the parent of
    $\pi_{m'+k+1}$, and for every $p$, $n_{j+k+1}=\suc{n_j}p$ iff
    $n_{m'+k+1}=\suc{n_{m'}}p$; or
  \item $n_{j+k}$ is a constant node and either $\pi_{m'+k}$ is the
    final position or $n_{j+k+1}\ne n_{m'+k+1}$.
  \end{enumerate}
\item A chain $\pi_{j_1},\dots,\pi_{j_{2p}}$ has the \emph{recurrence
  property} if for every $i\le p$ such that $n_{j_{2i-1}}$ is a variable
  node there are $j<j_{2i-1}$ and $k$ with $\iv{j+1}{j+k}$ and
  $\iv{m'+1}{m'+k}$ complementary, where $m'=j_{2i-1}$ and
  $j_{2i}=m'+k+1$.
\end{enumerate}
\end{definition}

\begin{axU}[S, Prop.~23]\label{U:between}
If $j<m$ and $\pi_i$ is the parent of both $\pi_j$ and $\pi_m$, then
there is $k$ with $j<k<m$ such that $\pi_j$ is the parent of $\pi_k$.
\end{axU}

\begin{axU}[S, Prop.~24]\label{U:compl}
If $\pi_i$ is the parent of $\pi_j$ and $\pi_m$, $j<m$, $n_j\ne n_m$ and
both carry the same variable, then there is $k$ such that
$\iv{j+1}{j+k}$ and $\iv{m+1}{m+k}$ are complementary.
\end{axU}

\begin{axU}[S, Fact~25]\label{U:end-child}
If $n_i$ is end then at most one position $\pi_j$, $j>i$, is a child of
$\pi_i$.
\end{axU}

\begin{axU}[S, Fact~26]\label{U:level1}
If $n_j,n_m$ ($j<m$) are level-$1$ nodes of $t$ carrying the same
variable, there is $k$ such that $\iv{j+1}{j+k}$ and $\iv{m+1}{m+k}$ are
complementary.
\end{axU}

\begin{axU}[S, Fact~27]\label{U:embedded}
Let $n,m$ carry the same variable, $m$ below $\suc np$, and $n_j=m$.
Then there are $i<j$ and $k$ with (1) $n_i=n$ and $\iv{i+1}{i+k}$,
$\iv{j+1}{j+k}$ complementary; (2) if $n,m$ are end and
$r_{i+1}=r_{i+k}$ then $n_{j+k+1}=\suc{n_j}p$.
\end{axU}

\begin{axU}[S, Prop.~28]\label{U:family}
Let $\pi_j$ be a position in $w_q$ and $\pi_{j_1},\dots,\pi_{j_{2p}}$ a
chain for $\pi_j$.  If $\pi_{j_{2k-1}},\pi_{j_{2m-1}}$ ($1\le k\le m\le p$)
are positions in $t$, then $n_{j_{2k-1}}$ and $n_{j_{2m-1}}$ belong to the
same family.
\end{axU}

\begin{axU}[S, Prop.~29, with the uniqueness shown in its proof]\label{U:maxchain}
Let $n_i$ be a variable or constant node in some $w_q$.  The chains
starting at $i$ with the recurrence property are totally ordered by
the prefix order and there is a longest one, of length $2r$ with
$r\le$ the depth of $w_q$.
\end{axU}

\begin{remark}
In Axiom~\ref{U:embedded}(2) the paper writes $[r_{i+1}]=[r_{i+k}]$,
i.e.\ the interval does not contribute.
Axiom~\ref{U:maxchain} is stated in the paper as ``there is a largest
$p$''; the definition of stages (Definition~\ref{def:partition}) needs
the chain itself to be unique, which is what the proof of [S, Prop.~29]
constructs (``we can uniquely extend the chain'').
\end{remark}

%======================================================================
\section{Two transformations}
%======================================================================

\begin{definition}[S, Section~6.2]
For a node $n$ of the $t$-component that is a variable or a constant
node and a term $s:\Ground$, $t[s/n]$ is $t$ with the subtree at $n$
replaced by the tree of $s$.  The \emph{size measure} of $t$ is
$\mu(t)=(|t|,\nu(t))$ ordered lexicographically, where $|t|$ is the
number of nodes (dummy lambdas included) and $\nu(t)$ the number of
variable nodes.  A solution of $E$ in $T_A(\CE)$ is \emph{smallest} if it
is $\mu$-minimal among such solutions.
\end{definition}

By Proposition~\ref{prop:CE} and well-foundedness of $\mu$, if $E$ has a
solution it has a smallest one.

\begin{definition}[$T_1$; S, Def.~38]
If $n$ is a variable node, or a constant node of positive arity, and no
play of $\Gm(t,E)$ has a position at $n$, then $t\,T_1\,t[b/n]$.
\end{definition}

\begin{proposition}[S, Fact~30]\label{prop:T1}
If $t$ solves $E$ and $t\,T_1\,t'$, then $t'$ solves $E$ and
$\mu(t')<\mu(t)$.
\end{proposition}
\begin{proof}
$t'$ is closed of type $A$ because the replaced subtree has type
$\Ground$.  By Lemma~\ref{lem:ancestors} no play visits a node below $n$
either.  Let $X$ be the set of nodes not below $n$; it is prefix closed
and the trees agree on it, so Lemma~\ref{lem:locality} gives
$\Gm(t',E)=\Gm(t,E)$, and $\forall$ loses it; by Axiom~\ref{G:char} $t'$
solves $E$.  If the replaced subtree has more than one node, $|t'|<|t|$;
otherwise $n$ is a variable of type $\Ground$ and $|t'|=|t|$,
$\nu(t')=\nu(t)-1$.
\end{proof}

\begin{definition}[End path; S, Def.~39]\label{def:endpath}
A downward path $\gamma=m_1,\dots,m_{2p}$ ($p\ge1$) of the $t$-component
is \emph{end} if
\begin{enumerate}
\item $m_1$ is a variable node (so the $m_{2i}$ are lambda nodes and the
  $m_{2i-1}$ are variable or constant nodes);
\item for $1<i\le p$, $m_{2i-1}$ is a variable node whose binder lies on
  $\gamma$;
\item every variable node strictly below $m_{2p}$ is bound by a node
  strictly above $m_1$ or strictly below $m_{2p}$.
\end{enumerate}
\end{definition}

\begin{remark}
The paper says ``below $m_{2p}$''.  It must be read strictly: if $m_{2p}$
bound a variable below it, that variable would become free after $T_2$.
In the paper's example the end path $(8),(11)$ satisfies the strict
reading.
\end{remark}

\begin{definition}[Redundancy; S, Def.~41]\label{def:redundant}
Let $\gamma$ be an end path and $\pi\in\Gm(t,E)$.  $\gamma$ is
\emph{redundant in} $\iv ik$ if either no $j_1$ with $i\le j_1<k$ has
$n_{j_1}=m_1$, or, for the least such $j_1$:
\begin{enumerate}[(a)]
\item there are $j\le k$ and a chain $\pi_{j_1},\dots,\pi_{j_{2p}}$ for
  $\pi_j$ with $n_{j_h}=m_h$ for all $h$ (so $n_j=m_{2p}$);
\item $\iv{j_1}j$ does not contribute;
\item $\gamma$ is redundant in $\iv jk$.
\end{enumerate}
(The definition is by induction on $k-i$, since $j>i$.)  $\gamma$ is
\emph{redundant in} $\Gm(t,E)$ if it is redundant in $\iv1l$ for every
play $\pi$ of length $l$.
\end{definition}

\begin{definition}[$T_2$; S, Def.~42]
If $\gamma$ is an end path redundant in $\Gm(t,E)$ and $t'$ is the
subterm at $\suc{m_{2p}}1$, then $t\,T_2\,t[t'/m_1]$.
\end{definition}

\begin{lemma}\label{lem:T2-term}
If $t\,T_2\,t''$ then $t''$ is a closed lnf of type $A$ with constants in
$\CE$ and $\mu(t'')<\mu(t)$.
\end{lemma}
\begin{proof}
Both replaced and replacing subtrees have type $\Ground$.  By
Definition~\ref{def:endpath}(3) every free variable of $t'$ is bound
strictly above $m_1$, so $t''$ is closed.  $t'$ is a proper subtree of
the replaced one, so $|t''|<|t|$.
\end{proof}

\begin{proposition}[S, Prop.~31]\label{prop:T2}
If $t$ solves $E$ and $t\,T_2\,t''$, then $t''$ solves $E$.
\end{proposition}
\begin{proof}[Proof structure]
Let $\pi\in\Gm(t,E)$.
\emph{Step 1 (decomposition; proved).} Unfolding
Definition~\ref{def:redundant} along $\iv1l$ gives
$\pi=\sigma_1\,\iv{i_1}{k_1}\,\sigma_2\cdots\iv{i_q}{k_q}\,\sigma_{q+1}$
where $\pi_{i_r}$ is the first position at $m_1$ after $\pi_{k_{r-1}}$,
$\pi_{k_r}$ is at $m_{2p}$, there is a chain following $\gamma$ from
$\pi_{i_r}$ to $\pi_{k_r}$, no $\iv{i_r}{k_r}$ contributes, and no
position of any $\sigma_r$ is at $m_1$.

\emph{Step 2 (no stray positions; sketched in the paper).} Every position
of a $\sigma_r$ that is below $m_1$ is below $m_{2p}$.  The paper argues
by contradiction with the earliest offending lambda position, using
Axioms~\ref{G:tables}, \ref{G:mono} and the end-path condition.

\emph{Step 3 (simulation).} Let $\rho$ map the node $\suc{m_{2p}}{1}w$ to
$m_1w$ and fix all nodes not below $m_1$.  Then $\pi'=\rho(\sigma_1)\cdots
\rho(\sigma_{q+1})$, with entries for binders on $\gamma$ removed from all
tables, is a play of $\Gm(t'',E)$; conversely every play of $\Gm(t'',E)$
arises in this way.  Since the removed intervals do not contribute,
$\pi$ and $\pi'$ end in the same final state, so $\forall$ loses
$\Gm(t'',E)$ and Axiom~\ref{G:char} gives the claim.
\end{proof}

\begin{gap}[Steps 2--3 of Proposition~\ref{prop:T2}]\label{gap:T2}
Needed: a \emph{simulation lemma} stating that after deleting the
intervals $\iv{i_r}{k_r}$, relocating nodes by $\rho$ and pruning the
tables, every move of $\Gm(t,E)$ in the $\sigma_r$ is a move of
$\Gm(t'',E)$, and conversely.  The delicate point is that tables in
$\sigma_{r+1}$ may hold entries created inside $\iv{i_r}{k_r}$ (jumps
from variables on $\gamma$ into $w$'s create entries pointing to
successors of nodes on $\gamma$); one must show that no later position
consults such an entry.  The paper's argument is a sketch and the
converse direction is ``almost identical''.  The statement is purely
game-level and could be added as an axiom of the U-layer.
\end{gap}

\begin{proposition}[S, Prop.~32]\label{prop:stacking}
Let $t$ be a smallest solution, and $m_1,\dots,m_p$ embedded end nodes of
$t$ carrying the same variable bound at the same node, with $m_{i+1}$
below $m_i$.  Then $p\le\depth(u)+1$.
\end{proposition}
\begin{proof}[Proof sketch, following the paper]
Let $m_p$ be below the $l$-th successor of $m_{p-1}$ and $m=\suc{m_p}l$.
Since $m_p$ is end, $m$ binds nothing, so $m_p,m$ is an end path
(Definition~\ref{def:endpath}).  If it were redundant in $\Gm(t,E)$, then
$T_2$ would apply and give a smaller solution
(Proposition~\ref{prop:T2}, Lemma~\ref{lem:T2-term}), a contradiction.
So in some play $\pi$ there is a position $\pi_j$ at $m_p$ whose
excursion $\iv{j+1}{j+k'}$ either contributes or does not return to $m$.
By Axiom~\ref{U:embedded}(2) applied to $(m_i,m_p)$ the second case
forces a contributing interval after an earlier $m_i$, and in the first
case Axiom~\ref{U:embedded}(1) gives intervals after positions at
$m_1,\dots,m_{p-1}$ complementary to (prefixes of) the excursion at
$m_p$.  These intervals contribute by Lemma~\ref{lem:corr-contrib} and
are pairwise disjoint, so Lemma~\ref{lem:disjoint} bounds their number by
$\depth(u)$.
\end{proof}

\begin{gap}[Proposition~\ref{prop:stacking}]\label{gap:stacking}
Two points are not argued in the paper: (i) that the complementary
intervals given by Axiom~\ref{U:embedded} for $m_1,\dots,m_{p-1}$ are
pairwise disjoint and ordered (plausible: $m_{i+1}$ is reached only
after the excursion from the position at $m_i$ has returned to a child
of that position); (ii) that the excursion at $m_p$ that witnesses
non-redundancy is the one to which Axiom~\ref{U:embedded} applies with
the correct length $k$.  The paper states the bound as $p\le\depth(u)$;
the counting above gives $p-1$ intervals, hence $p\le\depth(u)+1$, and
the paper's form needs that the excursion at $m_p$ itself contributes as
well, which is the first case.
\end{gap}

%======================================================================
\section{Partitioning plays}
%======================================================================

Throughout this section $t$ solves $E$ and $\pi\in\Gm(t,E)$ has length
$l$.

\begin{definition}[Stages; S, Def.~43]\label{def:partition}
Stage $1$ is $\iv12$.  Let stage $N$ be $\iv{a}{i-1}$ with $i-1<l$ and
$n_{i-1}$ a lambda node of $t$.  Then $n_i$ is the successor of
$n_{i-1}$ (Lemma~\ref{lem:pred}) and stage $N+1$ is defined by:
\begin{enumerate}
\item if $n_i$ is a constant node: $\iv i{i+1}$;
\item if $n_i$ is a variable node: $n_{i+1}$ is a lambda node of some
  $w_q$ (Axiom~\ref{G:alt}) and $n_{i+2}$ its successor.  Let
  $\pi_{j_1},\dots,\pi_{j_{2r}}$ ($j_1=i+2$) be the longest chain starting
  at $i+2$ with the recurrence property (Axiom~\ref{U:maxchain}).  Let
  $p$ be the least $p'\in\{1,\dots,r\}$ such that $\pi_{j_{2p'}+2}$ is a
  child of $\pi_i$, or $p=r$ if there is none.  Stage $N+1$ is
  $\iv i{i+3}$ if $p=0$ and $\iv i{j_{2p}+2}$ otherwise; the chain
  $\pi_{j_1},\dots,\pi_{j_{2p}}$ is \emph{associated} with stage $N+1$.
\end{enumerate}
\end{definition}

\begin{lemma}[Stages are well defined; S, Fact~33]\label{lem:partition}
The stages are consecutive, non-overlapping intervals covering $\iv1l$;
every non-final stage ends at a lambda node of $t$ and the last stage
ends at $\pi_l$.  The partition is unique.
\end{lemma}
\begin{proof}
Stage $1$ ends at node $1$, a lambda node of $t$.  Assume stage $N$ ends
at $\pi_{i-1}$, a lambda node of $t$.  Constant case: if $n_i$ is
$a:\Ground$, $\pi_{i+1}$ is final; if $n_i$ has positive arity,
$\pi_{i+1}$ is at a successor of $n_i$, a lambda node of $t$.  Variable
case, $p=0$: then $r=0$ (if $r\ge1$ then $p\ge1$).  If $n_{i+2}$ is a
variable node, $\pi_{i+3}$ is a jump into $t$ (Axiom~\ref{G:alt}).  It
cannot be a constant node of positive arity, since then
$\pi_{i+2},\pi_{i+3}$ would be a chain with the recurrence property
(Definition~\ref{def:parent}(2), constant case), so $r\ge1$.  If it is
$a:\Ground$, $\pi_{i+3}$ is final.  Variable case, $p\ge1$: if $p<r$ or
the cut condition holds at $p=r$, $\pi_{j_{2p}+2}$ is a child of $\pi_i$,
hence at a successor of $n_i$ (Axiom~\ref{G:parent}), a lambda node of
$t$.  Otherwise $p=r$; $\pi_{j_{2r}}$ is at a lambda node of $w_q$ and
$\pi_{j_{2r}+1}$ at its successor.  If that is a constant node of
positive arity, the chain could be extended by
$\pi_{j_{2r}+1},\pi_{j_{2r}+2}$ keeping the recurrence property,
contradicting maximality.  If it is $a:\Ground$, $\pi_{j_{2r}+2}$ is
final; if it is a variable node, $\pi_{j_{2r}+2}$ is a jump into $t$.
Every stage is non-empty and advances the index, and plays are finite,
so the stages cover $\pi$.  Uniqueness follows from
Axiom~\ref{U:maxchain}.
\end{proof}

\begin{remark}[On $\pi_{j_1}$]
By Definition~\ref{def:parent}(3) the odd members of a chain are
parents of lambda positions, so they are at variable, constant or $@$
nodes; in particular $\pi_{j_1}=\pi_{i+2}$ is at the successor of the
lambda node $n_{i+1}$ and is \emph{not} at a lambda node (it is an odd
position, Axiom~\ref{G:parity}(2)).  In [S, Example~14], stage
$S_5=\iv{15}{20}$ has associated chain $\pi_{17},\pi_{18}$ with $n_{17}$
the constant node $(32)$.  The lambda node of the stage is $n_{i+1}$,
which is not part of the associated chain.
\end{remark}

\begin{definition}[S, Def.~44]
Stage $N=\iv ij$ is \emph{final} if $r_j$ is final.  A non-final stage
\emph{contributes} if $\iv ij$ does, a final one if $\iv i{j-1}$ does
(Definition~\ref{def:contrib}).
\end{definition}

\begin{corollary}\label{cor:contrib-stages}
In every play at most $\depth(u)$ stages contribute; over the tree of
all plays at most $|u|-\br(u)$ distinct stages contribute.
\end{corollary}
\begin{proof}
Lemmas~\ref{lem:disjoint} and~\ref{lem:global-contrib}, together with
Lemma~\ref{lem:states}.
\end{proof}

\begin{proposition}[S, Prop.~34]\label{prop:stages}
Let $\iv ij$ be stage $N$ of $\pi$.
\begin{enumerate}
\item If $n_i$ is labelled $a:\Ground$, stage $N$ is final.
\item If $N$ is not final and $n_i$ is a constant node or a level-$1$
  node, then $\pi_j$ is a child of $\pi_i$.
\item If $N$ does not contribute and $n_i$ is embedded, $\pi_j$ is a
  child of $\pi_i$.
\item If $N$ is not final and $n_i$ is a variable node, $n_j$ is a
  successor of a node in the same family as $n_i$.
\end{enumerate}
\end{proposition}
\begin{proof}
(1) By Axiom~\ref{G:char} the $a$-move ends the play.

(2) Constant case: $j=i+1$ and the parent is $\pi_i$
(Definition~\ref{def:parent}(2)).  Level-$1$ case: $n_i$ is $x_q$, bound by
node $1$; by Axiom~\ref{G:mono}(2) applied to the chain from $\pi_1$ the
entry for node $1$ is always $((2,\dots,k+1),\varnothing)$, so
$n_{i+1}=q+1$ and $\theta_{i+1}(q+1)=((\suc{n_i}1,\dots),\theta_i)$.  The
constants met by the associated chain are first order (they are visited
with a state $f\,\vec r$ and $u$ is first order), so their successors
bind nothing; hence the first variable node $y_k$ met on the chain path
is bound by $q+1$.  By Axiom~\ref{G:mono}(2) applied to the chain
$\pi_1,\pi_{i+1},\pi_{i+2},\dots$, its position has parent $\pi_{i+1}$,
so the jump goes to $\suc{n_i}k$ and the resulting position is a child
of $\pi_i$ (Definition~\ref{def:parent}(2), variable case).  This is the
cut in Definition~\ref{def:partition}; if no variable is met, the stage
ends at $a:\Ground$ and is final.

(3) As in the paper, from Axiom~\ref{U:embedded}: the excursion is
complementary to an excursion from an earlier position at the
embedding node, which returned to a child.

(4) As in the paper, from Axiom~\ref{U:family}: extend the associated
chain to the chain starting at the successor of the root of $w_q$; the
position $\pi_{j-2}$ has a parent $\pi_{k'}$ among the chain positions
in $t$, so $n_j$ is a successor of $n_{k'-1}$, which is in the same
family as $n_i$.
\end{proof}

\begin{definition}[Comb, path; S, Def.~45]
Let $\iv ij$ be stage $N$.  The \emph{comb} $\comb(N)$ is $n_i$ together
with its successors.  The \emph{path} $\pth(N)$ is empty if $n_i$ is a
constant or $@$ node; if $n_i$ is a variable node with associated chain
$\pi_{j_1},\dots,\pi_{j_{2p}}$, it is
$n_{i+1},n_{j_1},\dots,n_{j_{2p}},n_{j-1}$, a downward path in some
$w_q$.  A stage is a \emph{constant/variable} stage, has \emph{level}
$k$, is \emph{embedded}, or belongs to a \emph{family}, according to
the root $n_i$ of its comb.
\end{definition}

\begin{definition}[Partition tree; S, Def.~46]\label{def:ptree}
The nodes are the stages of $\pi$.  For each non-final stage $N=\iv ij$,
there is one edge $N+1\ra\pi K$ where $K\le N$ is the largest stage whose
comb contains $n_j$ and such that $N\ra\pi^*M$ and $K\ra\pi^*M$ for some
$M\le N$.
\end{definition}

\begin{remark}\label{rem:ptree}
Every stage other than $1$ has exactly one outgoing edge, pointing to an
earlier stage, so $\ra\pi$ is a tree rooted at stage $1$ and every stage
reaches stage $1$.  Taking $M=1$, the second condition of
Definition~\ref{def:ptree} is therefore always satisfied, and the
definition reduces to ``the latest stage whose comb contains $n_j$''.
This is not obviously the intended relation: the same node can be the
comb root of several stages, visited with different tables
(Axiom~\ref{G:distinct}).  The reading that makes the reconstruction of
$t$ from combs correct is the following.
\end{remark}

\begin{definition}[Partition tree, proposed reading]\label{def:ptree2}
For a non-final stage $N=\iv ij$ let $\pi_h$ be the parent of $\pi_j$
(Axiom~\ref{G:parent}).  Put $N+1\ra\pi K$ where $K$ is the stage that
contains $\pi_h$.
\end{definition}

\begin{gap}[Partition tree]\label{gap:ptree}
Needed: (i) the parent $\pi_h$ in Definition~\ref{def:ptree2} is the
first position of its stage $K$ (so $n_j\in\comb(K)$); (ii) the two
definitions agree, or the proof of [S, Thm.~37] goes through with
Definition~\ref{def:ptree2}; (iii) \emph{coverage}: every node of $t$
visited by $\pi$ lies in the comb of some stage of $\pi$ (positions in
$t$ inside a stage lie in intervals complementary, hence corresponding,
to earlier ones, and induction on the index should give this).  With
(i)--(iii), gluing $\comb(N+1)$ below the lambda node of $\comb(K)$
along each edge reconstructs exactly the part of $t$ visited by $\pi$
[S, after Def.~46].
\end{gap}

\begin{lemma}[S, Fact~35]\label{lem:contrib-kinds}
If stage $N=\iv ij$ contributes then $N$ is a constant stage, or
$\pth(N)$ contains a constant node, or some $n_k$, $i<k<j$, is a
constant node of $t$ lying in the comb of an earlier stage $K$ of the
same family as $N$.
\end{lemma}
\begin{proof}[Proof sketch]
By Lemma~\ref{lem:states} the stage contains an $f$-move at some $n_k$.
If $n_k$ is in $w_q$ it lies on the chain path, i.e.\ on $\pth(N)$.  If
$n_k=n_i$ the stage is a constant stage.  Otherwise $n_k$ lies in $t$
inside a recurrence excursion $\iv{m'+1}{m'+k'}$, which corresponds to an
earlier excursion; by Gap~\ref{gap:ptree}(iii) $n_k$ lies in an earlier
comb, and the family claim is Axiom~\ref{U:family}.
\end{proof}

%======================================================================
\section{Unfolding and the small solution theorem}
%======================================================================

\begin{definition}[Unfolding; S, Def.~47]\label{def:unfold}
Let $N_1<\dots<N_k$ be stages of $\pi$ with combs $c_i$ and paths
$\gamma_i$ such that
\begin{enumerate}[(a)]
\item they belong to the same family and do not contribute;
\item $c_i\ne c_j$ and $\gamma_i\ne\gamma_j$ for $i<j$;
\item for every $j\le k$ (except $j=k$ if $N_k$ is final) there is
  $i\le j$ with $N_j+1\ra\pi N_i$;
\item for every $1<j\le k$ there is $i<j$ such that the root of
  $\gamma_j$ is a successor of the last node of $\gamma_i$.
\end{enumerate}
The \emph{enlarged comb} $(c_1\dots c_j)$ is the tree built from copies
of $c_1,\dots,c_j$: $(c_1)=c_1$, and $(c_1\dots c_{j+1})$ is obtained by
placing the root of a fresh copy of $c_{j+1}$ below the copy of the
lambda leaf at which $N_j$ finishes (by (c) this leaf belongs to the
copy of some $c_i$, $i\le j$).  Variables in the copy of $c_{j+1}$ are
bound by the copies of their binders inside $(c_1\dots c_j)$ when those
binders belong to some $c_i$; otherwise they keep their binders.  The
\emph{enlarged tree} $(\gamma_1\dots\gamma_j)$ is the union of the paths
$\gamma_1,\dots,\gamma_j$, a subtree of $w_q$ by (d).
\end{definition}

\begin{definition}[Unfolded partition tree]\label{def:unfolded-tree}
Unfolding $N_1$ into $N_2,\dots,N_k$ replaces, for $j\ge2$, the stage
$N_j$ by a new stage $N'_j$ with comb $(c_1\dots c_j)$ and the outgoing
edge of $N_j$, and redirects every edge $N_j+1\ra\pi N_i$ given by (c) to
$N_j+1\ra{} N'_j$, attached at the copy of the leaf of $c_i$ inside
$(c_1\dots c_j)$.  The \emph{unfolded term} is obtained by gluing the
combs of the resulting tree along its edges (with $\lambda\vec y.b$ below
lambda leaves that receive no edge).
\end{definition}

\begin{example}[S, Example~14, Figs.~12--15]
$S_3$ is unfolded into $S_6$ and into $S_8$.  $c_3$ is $(4)x_1$ with
successors $(5)$ and $(7)\lambda s_1s_2$; $c_6=(16)s_2$.  In $(c_3c_6)$ the
copy of $s_2$ is placed below the copy of $(7)$ and is bound by it; the
play of $(S_3S_6)$ follows $\gamma_3=(17),(18)$, returns to $(7)$, then
follows $\gamma_6=(21),(22)$ and returns to $(5)$, where $S_6$ originally
ended.  $S_7=S_6+1$, which pointed to $S_3$, now points to $S'_6$.
\end{example}

\begin{lemma}[S, Fact~36, local part]\label{lem:unfold-local}
In the game for the unfolded term, the play from the root of
$(c_1\dots c_j)$ follows the enlarged tree $(\gamma_1\dots\gamma_j)$, does
not contribute, and finishes at the copy of node $n$ of $c_i$ iff
$N_j$ finishes at $n$ of $c_i$.
\end{lemma}
\begin{proof}[Proof sketch]
Induction on $j$.  Inside $(c_1\dots c_j)$ the binders of the copies are
local, and by (d) the jumps from the copy of $c_{j}$ go into
$\gamma_j$ with the table entries created at the copy of $c_{i}$, which
correspond to the entries used in the original stage $N_j$ (by
Axiom~\ref{G:agree}, entries agree on binders above).  The states do
not change, since no stage $N_i$ contributes.
\end{proof}

\begin{gap}[Soundness of unfolding]\label{gap:unfold}
The paper does not define the game for the unfolded term (``we do not
provide the game to save notation'') and does not prove that the
unfolded term solves $E$.  Needed: (i) a precise statement of
Lemma~\ref{lem:unfold-local} for the new game, including the tables at
which the new stage ends; (ii) \emph{global soundness}: every play of
the game of the unfolded term decomposes into stages that are either
original stages or new stages $N'_j$, with the same final state, so
$\forall$ loses (Axiom~\ref{G:char}); (iii) \emph{consistency across
plays}: the term is shared by up to $\br(u)$ plays, and the partition
tree is per play.  An unfolding performed for one play changes the
subterm used by the others; the paper handles this with
\emph{separators} (stages after which two plays sharing all earlier
combs finish at different nodes) only in the final count.
\end{gap}

\begin{definition}[c-embedded; S, Def.~48]
An enlarged comb $(c_1\dots c_k)$ is \emph{c-embedded} if its only free
variable occurrence is the head of $c_1$ and some enlarged comb
$(c'_1\dots c'_k)$ above it in the partition tree is $\alpha$-equivalent
to it.
\end{definition}

\begin{theorem}[S, Thm.~37]\label{thm:small}
Let $E$ have order $2p$ or $2p+1$ and let $q$ bound the arity of $E$ and
of the constants of $u$.  If $E$ has a (canonical) solution then it has
a (canonical) solution of depth at most $4p^2q^{2p-2}|u|$.
\end{theorem}

\begin{proof}[Proof architecture]
Let $t$ be a smallest solution.  We list the steps of the paper with
their status.
\begin{enumerate}[{C}1.]
\item \emph{(Proved.)} No $T_1$ or $T_2$ applies to $t$
  (Propositions~\ref{prop:T1}, \ref{prop:T2} and
  Lemma~\ref{lem:T2-term}; the latter depends on Gap~\ref{gap:T2}).
  Hence every variable node and every constant node of positive arity
  is visited by some play.
\item \emph{(Proved.)} Every node of $t$ has level at most $p$
  (Lemma~\ref{lem:levels}); branching and binder length are at most
  $q$; on every path at most $\depth(u)$ constant nodes have positive
  arity (Lemma~\ref{lem:constants}).
\item \emph{(Proved.)} If a binder binds $(k+1)q$ nodes on one path, then
  at least $k$ of them are embedded nodes carrying the same variable
  (pigeonhole: the binder binds at most $q$ variables, so one variable
  occurs at least $k+1$ times; all but the highest are embedded).
\item \emph{(Proved modulo Gap~\ref{gap:stacking}.)} Among embedded end
  nodes carrying the same variable on one path there are at most
  $\depth(u)+1$ (Proposition~\ref{prop:stacking}).
\item \emph{(Gap~\ref{gap:unfold}.)} If $k>q$ embedded non-end nodes of
  level $j$ carry the same variable on a path, a stage $N_1$ at one of
  them can be unfolded so that the enlarged comb becomes c-embedded; the
  resulting term is again a solution.  Unfolding is applied exhaustively
  from stages that have no contributing descendant under binding,
  starting at level $1$ or embedded stages; the depth of an enlarged
  comb is at most $2pq^{p-1}$ (``a level-$1$ node is unfolded into $q$
  level-$2$ nodes, which are unfolded into $q^2$ level-$3$ nodes, and so
  on'').  Termination of the combined procedure (unfolding increases
  size, $T_1,T_2$ decrease it) is not argued.
\item \emph{(Gap.)} After exhaustive unfolding, $T_2$ applies to every
  enlarged comb whose play runs from its root to a leaf along a
  redundant end path.  The stages to which $T_2$ does not apply are the
  contributing ones (at most $|u|$ over all plays,
  Corollary~\ref{cor:contrib-stages}) and the later non-embedded stages
  starting from descendants under binding of their roots (at most
  $pq^{p-1}$ per contributing stage, not argued), plus separators (at
  most $pq^{p-1}|u|$, not argued).  Multiplying by the depth
  $2pq^{p-1}$ of an enlarged comb and by $2$ for the lambda layers gives
  the bound $4p^2q^{2p-2}|u|$.
\end{enumerate}
The theorem follows from C1--C6.  C1--C3 are proved here, C4 depends on
Gap~\ref{gap:stacking}, and C5--C6 constitute Gap~\ref{gap:count}.
For canonical solutions, note that $T_1$, $T_2$ and unfolding never
touch the head $z$ of $\lambda z.z\,\vec t$ nor create free occurrences
of $z$ in the $t_i$ ($z$ is level $1$, and its node is the first
non-lambda node visited), so canonicity is preserved.
\end{proof}

\begin{gap}[Counting in Theorem~\ref{thm:small}]\label{gap:count}
Needed, with Definition~\ref{def:ptree2} and after Gap~\ref{gap:unfold}:
(i) a measure that decreases under the combined procedure
``unfold exhaustively, then apply $T_1,T_2$ exhaustively''; (ii) for each
level $j$, the bound ``a path containing $(k+1)q^j$ level-$j$ nodes
contains $k$ embedded nodes with the same variable'' (the paper's
induction on $j$ is not spelled out: level-$j$ nodes on a path can be
bound by different binders); (iii) the bound on non-embedded
descendants under binding of a contributing comb; (iv) the bound on
separators across plays.  Each of these is a statement about the
partition trees of the plays of a fixed solution, so after
Gap~\ref{gap:unfold} they are finite combinatorics.
\end{gap}

\begin{corollary}[Decidability]\label{cor:decidable}
Assume Theorem~\ref{thm:small}.  Higher-order matching (under (R1),
(R2)) is decidable.
\end{corollary}
\begin{proof}
By Proposition~\ref{prop:canonical} it suffices to decide whether the
associated interpolation equation $E$ has a canonical solution, and by
Proposition~\ref{prop:CE} we may look in $T_A(\CE)$.  Let
$D=4p^2q^{2p-2}|u|$.  Up to $\alpha$-equivalence the lnf terms of type $A$
over $\CE$ of depth $\le D$ form a finite, effectively enumerable set:
branching and binder lengths are bounded by $q$, and on each path at most
$D$ binders occur, so variables can be named from a fixed set of
$qD$ names (each binder at depth $d$ uses names
$v_{d,1},\dots,v_{d,q}$).  For each such canonical term $t$ compute the
$\beta$-normal form of $t\,\vec w$ (Fact~\ref{L:basic}) and compare with
$u$ (Lemma~\ref{lem:reduces}).
\end{proof}

%======================================================================
\section{Removing the restrictions}\label{sec:restrictions}
%======================================================================

[S, Section~6.5] sketches two extensions; neither is formalized here.
\begin{itemize}
\item \emph{Several base types.} Nodes carry their result type; a play
  changes result type only at constant nodes, so at most $\depth(u)$
  times; $T_1$ uses a constant $b_B$ for each base type $B$, and $T_2$
  and unfolding are restricted to paths of constant result type.  The
  effect on Theorem~\ref{thm:small} is to replace $|u|$ by $|u|^2$.
\item \emph{Goal with bound variables.} Look-up tables get entries for
  binders below constants, mapping their variables to fresh constants
  $c_i$; the constant and variable moves are amended [S, Fig.~16].  The
  paper states that Theorem~13 holds for the new game, i.e.\ the
  G-axioms must be re-postulated for it, and that the transformations
  and unfolding are unchanged.
\end{itemize}

%======================================================================
\section{Summary of dependencies and gaps}\label{sec:gaps}
%======================================================================

\begin{center}
\begin{tabular}{lll}
Result here & Paper & Depends on\\\hline
Lemmas~\ref{lem:pred}--\ref{lem:ancestors} & after Def.~23 & G\ref{G:parent}, moves\\
Lemma~\ref{lem:locality} & (implicit) & moves only\\
Lemmas~\ref{lem:states}--\ref{lem:global-contrib} & (implicit) & G\ref{G:char}, G\ref{G:total}, G\ref{G:parity}\\
Prop.~\ref{prop:T1} & Fact~30 & Lemma~\ref{lem:locality}, G\ref{G:char}\\
Prop.~\ref{prop:T2} & Prop.~31 & Gap~\ref{gap:T2}\\
Prop.~\ref{prop:stacking} & Prop.~32 & U\ref{U:embedded}, Gap~\ref{gap:stacking}\\
Lemma~\ref{lem:partition} & Def.~43, Fact~33 & G\ref{G:alt}, U\ref{U:maxchain}\\
Prop.~\ref{prop:stages} & Prop.~34 & G\ref{G:mono}, U\ref{U:embedded}, U\ref{U:family}\\
Lemma~\ref{lem:contrib-kinds} & Fact~35 & Gap~\ref{gap:ptree}\\
Thm.~\ref{thm:small} & Thm.~37 & Gaps~\ref{gap:T2}--\ref{gap:count}\\
Cor.~\ref{cor:decidable} & Section~6 & Thm.~\ref{thm:small}
\end{tabular}
\end{center}

\paragraph{Gaps.}
\begin{enumerate}
\item Gap~\ref{gap:T2}: simulation lemma for $T_2$ (game-level).
\item Gap~\ref{gap:stacking}: disjointness and counting in Prop.~32.
\item Gap~\ref{gap:ptree}: the definition of the partition tree is
  degenerate as written (Remark~\ref{rem:ptree}); proposed reading via
  parents, plus coverage of $t$ by combs.
\item Gap~\ref{gap:unfold}: the game of an unfolded term, soundness of
  unfolding, consistency across plays.  This is the main gap.
\item Gap~\ref{gap:count}: termination of the procedure and the counting
  that yields $4p^2q^{2p-2}|u|$.
\end{enumerate}
If only Section~4 is taken as axiomatic, Axioms U\ref{U:between}--U\ref{U:maxchain}
become additional obligations; the paper's proofs of them use the move
rules, Axioms G\ref{G:parent}--G\ref{G:alt} and case analysis on node
labels.

\paragraph{Corrections and readings adopted.}
\begin{itemize}
\item End paths: ``below $m_{2p}$'' read as strictly below
  (Definition~\ref{def:endpath}).
\item Contribution extended to intervals not ending at a lambda node
  (Definition~\ref{def:contrib}), as needed for final stages.
\item Stage selection uses the cut ``first $p$ with
  $\pi_{j_{2p}+2}$ a child of $\pi_i$'' from the text after [S, Def.~43],
  which is more precise than the condition ``$\iv{j_1}{j_{2p}}$ contains
  no child of $\pi_i$'' in the definition itself.
\item Theorem~\ref{thm:small}: $q$ must also bound the arities of the
  constants of $u$, since constant nodes of $t$ branch according to
  them.
\item Prop.~\ref{prop:stacking}: the argument gives $p\le\depth(u)+1$;
  the paper's $p\le\depth(u)$ needs the extra observation in
  Gap~\ref{gap:stacking}.
\end{itemize}

\end{document}

%%% Local Variables:
%%% mode: latex
%%% TeX-master: t
%%% End:
